\chapter{Phases Explained}\label{ch:phases-explained}

\index{phase system|(}

Every material in nature exists in a state.
Water flows as liquid, hardens as ice, disperses as steam.
The state determines what you can \emph{do} with the material---you can pour water but not ice, you can stack ice but not steam.
Lattice borrows this intuition and applies it to data: every value in the language carries a \emph{phase} that governs whether it can be changed, shared, or locked down permanently.

If you have been following along from \cref{ch:variables-and-phase}, you already know the basics: \lstinline|flux| for mutable, \lstinline|fix| for immutable, and \lstinline|let| for ``figure it out for me.''
This chapter goes deeper.
We will explore the full philosophy behind the phase system, examine each phase keyword in detail, master the trio of \lstinline|freeze()|, \lstinline|thaw()|, and \lstinline|clone()|, learn when and why to freeze values, and understand the error messages the phase checker produces when something goes wrong.

\section{The Philosophy: Mutability as a Material Property}\label{sec:mutability-as-material}

\index{mutability}

Most languages treat mutability as a \emph{declaration modifier}.
In Rust, you write \lstinline|let mut|.
In JavaScript, you choose between \lstinline|let| and \lstinline|const|.
In both cases, the declaration tells the \emph{binding} whether it can be reassigned, but the underlying data may or may not follow suit.

Lattice takes a different stance.
Mutability is not a property of the \emph{binding}---it is a property of the \emph{value itself}.
When you declare a variable with \lstinline|flux|, the value stored inside is tagged as \emph{fluid}: it lives in the mutable heap, the garbage collector watches over it, and you can modify it freely.
When you declare with \lstinline|fix|, the value is tagged as \emph{crystal}: it is hardened, immutable, and may be stored in a memory arena where it can never be altered.

\begin{defnbox}{Phase}
A \emph{phase}\index{phase!definition} is a runtime tag on every Lattice value that describes its mutability state.
There are four phase tags in the runtime:
\begin{itemize}
  \item \textbf{Fluid} (\lstinline|VTAG_FLUID|) --- mutable, alive, in motion.
  \item \textbf{Crystal} (\lstinline|VTAG_CRYSTAL|) --- immutable, hardened, permanent.
  \item \textbf{Unphased} (\lstinline|VTAG_UNPHASED|) --- newly created, not yet committed to either side.
  \item \textbf{Sublimated} (\lstinline|VTAG_SUBLIMATED|) --- permanently immutable; cannot be thawed back.
\end{itemize}
\end{defnbox}

This design has a profound consequence: you can take a fluid value and \emph{freeze} it into a crystal.
You can take a crystal and \emph{thaw} it back into a fluid.
You can \emph{clone} a value to get an independent copy without changing the original's phase.
The phase is the material property of the data itself, and Lattice gives you the tools to change that material state deliberately.

Why does this matter?
Consider a configuration object.
While your application is starting up, the configuration needs to be mutable---you are loading defaults, overriding with environment variables, parsing command-line flags.
But once startup is complete, the configuration should be frozen solid: no accidental mutation, no race conditions, no surprises.
In most languages, you enforce this with discipline and documentation.
In Lattice, you enforce it with a single call to \lstinline|freeze()|.

\begin{lstlisting}[language=Lattice, caption={Freezing a configuration after setup}]
flux config = Map::new()
config.set("host", "0.0.0.0")
config.set("port", "8080")
config.set("debug", "false")

// Startup complete. Lock it down.
fix app_config = freeze(config)
print(phase_of(app_config))  // "crystal"
\end{lstlisting}

After the freeze, \lstinline|app_config| is crystal.
Any attempt to call \lstinline|app_config.set("host", "evil.com")| will fail with a phase error at runtime.
The data is safe---not because of a linting rule or a code review comment, but because the material itself has hardened.

\section{\texttt{flux} (Fluid) --- Mutable, Alive, in Motion}\label{sec:flux-fluid}

\index{flux@\texttt{flux}|(}
\index{fluid phase}

The \lstinline|flux| keyword declares a binding whose value is in the \emph{fluid} phase.
Fluid values can be reassigned, mutated in place, grown, shrunk, and generally treated as living, changing data.

\begin{lstlisting}[language=Lattice, caption={Basic \texttt{flux} usage}]
flux counter = 0
counter = counter + 1
counter += 1
print(counter)  // 2
\end{lstlisting}

When you create a value with \lstinline|flux|, the runtime tags it as \lstinline|VTAG_FLUID|.
Under the hood (in \filename{include/value.h}), every \lstinline|LatValue| struct carries a \lstinline|phase| field of type \lstinline|PhaseTag|.
A fluid value's phase tag is set to \lstinline|VTAG_FLUID|, and it lives on the \emph{FluidHeap}---the garbage-collected portion of memory managed by the mark-sweep collector.

\subsection{Fluid Collections}

Fluid values are not limited to scalars.
Arrays, maps, sets, and structs can all be fluid:

\begin{lstlisting}[language=Lattice, caption={Fluid collections}]
flux temperatures = [18.5, 19.2, 20.1]
temperatures.push(21.0)
temperatures.push(22.3)
print(temperatures)  // [18.5, 19.2, 20.1, 21.0, 22.3]

flux users = Map::new()
users.set("alice", 42)
users.set("bob", 37)
print(users.len())  // 2
\end{lstlisting}

When a compound value like an array is fluid, all of its elements inherit the fluid phase.
You can modify the array and its contents freely.

\subsection{Fluid Reassignment}

Because \lstinline|flux| values are mutable bindings, you can reassign them entirely:

\begin{lstlisting}[language=Lattice, caption={Reassigning a \texttt{flux} binding}]
flux greeting = "hello"
print(phase_of(greeting))  // "fluid"

greeting = "bonjour"
print(greeting)  // "bonjour"

// Even change the type entirely
greeting = 42
print(greeting)  // 42
\end{lstlisting}

Lattice is dynamically typed, so a \lstinline|flux| binding can hold any value.
The phase stays fluid unless you explicitly transition it.

\begin{warnbox}{Fluid Does Not Mean Unprotected}
A fluid value is mutable, but it is still tracked by the runtime.
The garbage collector manages its memory, the phase checker watches for incorrect transitions, and in strict mode, the compiler verifies that you do not accidentally pass a fluid value where a crystal is expected.
\end{warnbox}

\index{flux@\texttt{flux}|)}

\section{\texttt{fix} (Crystal) --- Immutable, Hardened, Permanent}\label{sec:fix-crystal}

\index{fix@\texttt{fix}|(}
\index{crystal phase}

The \lstinline|fix| keyword declares a binding whose value is in the \emph{crystal} phase.
Crystal values are immutable: they cannot be reassigned, their internal state cannot be modified, and any attempt to mutate them produces a runtime error.

\begin{lstlisting}[language=Lattice, caption={Basic \texttt{fix} usage}]
fix pi = 3.14159
fix greeting = "Hello, Lattice!"
fix primes = freeze([2, 3, 5, 7, 11])

print(pi)        // 3.14159
print(greeting)  // Hello, Lattice!
print(primes)    // [2, 3, 5, 7, 11]
\end{lstlisting}

\begin{notebox}{Scalars with \texttt{fix}}
Scalar values (integers, floats, booleans, strings) bound with \lstinline|fix| are automatically crystal.
You do not need to call \lstinline|freeze()| on them---the runtime handles the phase assignment.
For compound values like arrays, maps, and structs, you typically call \lstinline|freeze()| explicitly to ensure the deep crystallization (more on this in \cref{sec:freeze-thaw-clone}).
\end{notebox}

\subsection{Why Crystal Is Not Just ``const''}

In many languages, \lstinline|const| merely prevents \emph{rebinding}---you cannot assign a new value to the variable, but you can still modify the data it points to.
JavaScript's \lstinline|const| is the canonical example:

\begin{lstlisting}[language=Lattice, caption={The problem with shallow const (JavaScript-style)}]
// In JavaScript:
// const arr = [1, 2, 3];
// arr.push(4);  // This works! const doesn't protect contents.

// In Lattice:
fix arr = freeze([1, 2, 3])
// arr.push(4)  // Runtime error: cannot mutate crystal value
\end{lstlisting}

Lattice's crystal phase is \emph{deep}.
When you freeze an array, every element becomes crystal.
When you freeze a struct, every field becomes crystal.
When you freeze a map, every key-value pair becomes crystal.
The immutability propagates through the entire object graph.

Under the hood, the \lstinline|value_freeze()| function in \filename{src/value.c} calls \lstinline|set_phase_recursive()|, which walks the entire value tree and sets every nested value's phase to \lstinline|VTAG_CRYSTAL|:

\begin{lstlisting}[language=Lattice, caption={Deep freeze in action}]
flux nested = [[1, 2], [3, 4], [5, 6]]
fix frozen = freeze(nested)

print(phase_of(frozen))  // "crystal"
// Every inner array is also crystal---you cannot push
// into frozen[0] any more than you can push into frozen itself.
\end{lstlisting}

\subsection{Crystal Values and Memory}

Crystal values may be relocated from the FluidHeap into a \emph{CrystalRegion}---an arena-backed memory area that provides cache locality and O(1) bulk deallocation.
We explore this memory architecture in detail in \cref{ch:memory-arenas}.
For now, the key insight is that crystal values are not merely ``values you are not allowed to change.''
They are values that \emph{the runtime treats differently at the memory level}, storing them in a region optimized for immutable data.

\index{fix@\texttt{fix}|)}

\section{\texttt{let} --- Phase Inference}\label{sec:let-inference}

\index{let@\texttt{let}|(}
\index{phase inference}

Not every binding needs an explicit phase declaration.
The \lstinline|let| keyword tells Lattice: ``I do not care about the phase right now---infer it from context.''

\begin{lstlisting}[language=Lattice, caption={\texttt{let} bindings with inferred phase}]
let name = "Lattice"
let count = 42
let items = [1, 2, 3]

print(phase_of(name))   // "unphased"
print(phase_of(count))  // "unphased"
print(phase_of(items))  // "unphased"
\end{lstlisting}

When you use \lstinline|let|, the value starts in the \lstinline|VTAG_UNPHASED| state.
Unphased values are not locked into either fluid or crystal.
They behave as fluid in practice---you can reassign and mutate them---but they do not carry the explicit guarantee of either phase.

\begin{defnbox}{Unphased}
\index{unphased}
An \emph{unphased} value (\lstinline|VTAG_UNPHASED|) has no explicit phase commitment.
It can be treated as mutable and can transition to either fluid or crystal through explicit operations.
In casual mode (the default), \lstinline|let| bindings are unphased.
In strict mode, \lstinline|let| is \emph{forbidden}---you must choose \lstinline|flux| or \lstinline|fix| explicitly.
\end{defnbox}

The phase checker in \filename{src/phase\_check.c} handles \lstinline|let| differently depending on the execution mode.
In casual mode (\lstinline|MODE_CASUAL|), the effective phase of a \lstinline|let| binding is inferred from the initializer expression:

\begin{lstlisting}[language=Lattice, caption={Phase inference with \texttt{let}}]
let frozen_data = freeze([1, 2, 3])
// frozen_data's effective phase is crystal (inferred from freeze())

let thawed_copy = thaw(frozen_data)
// thawed_copy's effective phase is fluid (inferred from thaw())
\end{lstlisting}

\begin{tipbox}{When to Use \texttt{let}}
Use \lstinline|let| for quick prototyping and REPL exploration.
When writing production code, prefer \lstinline|flux| or \lstinline|fix| to make your intent explicit.
In strict mode, you are \emph{required} to make this choice---the compiler will refuse \lstinline|let| bindings entirely.
\end{tipbox}

\subsection{The Relationship Between \texttt{let} and Strict Mode}

In strict mode (\lstinline|#mode strict|), the phase checker rejects \lstinline|let| bindings with a clear error:

\begin{lstlisting}[language=Lattice, caption={\texttt{let} rejected in strict mode}]
#mode strict

let name = "Lattice"
// Error: strict mode: use 'flux' or 'fix' instead of 'let'
//        for binding 'name'
\end{lstlisting}

This is by design.
Strict mode exists for codebases where phase discipline matters---library code, concurrent systems, safety-critical paths.
We cover strict mode in full detail in \cref{ch:strict-mode}.

\index{let@\texttt{let}|)}

\section{\texttt{freeze()}, \texttt{thaw()}, \texttt{clone()} in Depth}\label{sec:freeze-thaw-clone}

\index{freeze@\texttt{freeze()}|(}
\index{thaw@\texttt{thaw()}|(}
\index{clone@\texttt{clone()}|(}

These three built-in functions are the tools you use to change a value's phase or create independent copies.
Think of them as the laboratory equipment for Lattice's material science.

\subsection{\texttt{freeze()} --- Crystallization}\label{subsec:freeze}

\lstinline|freeze()| takes a value and returns a crystal copy.
The original fluid value is unaffected (though it may be freed by garbage collection if no longer referenced), and the returned value is deeply immutable.

\begin{lstlisting}[language=Lattice, caption={The \texttt{freeze()} function}]
flux temperatures = [20.1, 21.5, 19.8, 22.0]
fix snapshot = freeze(temperatures)

print(phase_of(temperatures))  // "fluid"
print(phase_of(snapshot))      // "crystal"

// The fluid value is still alive and mutable
temperatures.push(23.1)
print(temperatures)  // [20.1, 21.5, 19.8, 22.0, 23.1]

// The crystal snapshot is frozen in time
print(snapshot)  // [20.1, 21.5, 19.8, 22.0]
\end{lstlisting}

Under the hood, \lstinline|freeze()| does two things:

\begin{enumerate}
  \item \textbf{Deep clones} the value into a new \lstinline|CrystalRegion| (arena-backed memory). This is performed by \lstinline|value_deep_clone()| in \filename{src/value.c}, which recursively copies every nested value.
  \item \textbf{Sets the phase tag} to \lstinline|VTAG_CRYSTAL| on every value in the tree via the \lstinline|set_phase_recursive()| function.
\end{enumerate}

The deep clone ensures that the crystal value has no pointers back into the FluidHeap.
This is a critical invariant: crystal values must be completely independent of the garbage-collected heap so that the GC can skip them during mark-sweep cycles.

\begin{notebox}{Freeze Is Not Free}
\lstinline|freeze()| performs a deep copy.
For large nested structures, this has a real cost in both time and memory.
However, this cost buys you a powerful guarantee: the frozen value is truly independent and will never be mutated, even if the original value continues to change.
\end{notebox}

\subsubsection{Freezing with Contracts}

\lstinline|freeze()| can take an optional contract closure that validates the value before crystallization:

\begin{lstlisting}[language=Lattice, caption={Freeze with contract validation}]
flux score = 85

// freeze with a validation contract
fix validated = freeze(score, |val: any| {
    if val < 0 { return "score must be non-negative" }
    if val > 100 { return "score must not exceed 100" }
    return nil
})

print(validated)  // 85
\end{lstlisting}

If the contract returns a non-nil value, the freeze fails and the error message is propagated.
This pattern lets you enforce invariants at the moment of crystallization---the value only hardens if it passes validation.

\subsection{\texttt{thaw()} --- Melting a Crystal}\label{subsec:thaw}

\lstinline|thaw()| is the inverse of \lstinline|freeze()|.
It takes a crystal value and returns a \emph{new} fluid copy.
The original crystal value remains unchanged and immutable.

\begin{lstlisting}[language=Lattice, caption={The \texttt{thaw()} function}]
fix frozen_list = freeze([10, 20, 30])
print(phase_of(frozen_list))  // "crystal"

flux thawed = thaw(frozen_list)
print(phase_of(thawed))  // "fluid"

// Now we can modify the thawed copy
thawed.push(40)
print(thawed)       // [10, 20, 30, 40]
print(frozen_list)  // [10, 20, 30]  -- unchanged
\end{lstlisting}

The implementation in \filename{src/value.c} reveals an important detail: \lstinline|value_thaw()| first calls \lstinline|value_deep_clone()| to create an independent copy, then calls \lstinline|set_phase_recursive()| to set the phase to \lstinline|VTAG_FLUID|.
This means \lstinline|thaw()| always creates a new copy---it never modifies the original crystal value in place.

\begin{warnbox}{Thaw Creates a Copy}
A common misconception is that \lstinline|thaw()| ``unfreezes'' the original value.
It does not.
\lstinline|thaw()| creates a brand-new fluid copy, leaving the original crystal intact.
If you need to modify a frozen value, thaw it, modify the copy, and optionally re-freeze it.
\end{warnbox}

\subsubsection{Thawing Refs}

When \lstinline|thaw()| operates on a \lstinline|Ref| value (a reference-counted shared wrapper), it breaks the sharing.
The thawed result is a new \lstinline|Ref| with a deep-cloned inner value, independent of any other references to the original.
This is documented in \filename{src/value.c} in the \lstinline|value_thaw()| function:

\begin{lstlisting}[language=Lattice, caption={Thawing a Ref breaks sharing}]
fix shared_ref = freeze(Ref(42))
flux independent = thaw(shared_ref)
// independent is a new Ref, no longer sharing data with shared_ref
\end{lstlisting}

\subsection{\texttt{clone()} --- Independent Copy}\label{subsec:clone}

\lstinline|clone()| creates a deep copy of a value \emph{without changing its phase}.
A fluid value cloned stays fluid.
A crystal value cloned stays crystal.
The result is a structurally identical but completely independent copy.

\begin{lstlisting}[language=Lattice, caption={The \texttt{clone()} function}]
flux original = [1, 2, 3]
flux copy = clone(original)

copy.push(4)
print(original)  // [1, 2, 3]  -- unaffected
print(copy)      // [1, 2, 3, 4]

print(phase_of(original))  // "fluid"
print(phase_of(copy))      // "fluid"  -- same phase as original
\end{lstlisting}

Clone is especially useful when you want to pass a value to a function that might modify it, but you want to preserve the original:

\begin{lstlisting}[language=Lattice, caption={Cloning for safe function calls}]
fn process_data(data: any) {
    // This function modifies its argument
    data.push(0)
    return data
}

flux sensor_readings = [45.2, 47.8, 46.1]
flux processed = process_data(clone(sensor_readings))

print(sensor_readings)  // [45.2, 47.8, 46.1]  -- safe
print(processed)        // [45.2, 47.8, 46.1, 0]
\end{lstlisting}

\begin{tipbox}{When to Use Each}
\begin{itemize}
  \item Use \lstinline|freeze()| when you want to lock a value down for safety, sharing, or performance.
  \item Use \lstinline|thaw()| when you need to modify data that is currently frozen.
  \item Use \lstinline|clone()| when you want an independent copy without changing the phase.
\end{itemize}
\end{tipbox}

\index{freeze@\texttt{freeze()}|)}
\index{thaw@\texttt{thaw()}|)}
\index{clone@\texttt{clone()}|)}

\section{Advanced Phase Operations}\label{sec:advanced-phase-ops}

Beyond the fundamental trio, Lattice offers several advanced phase operations for fine-grained control over the crystallization lifecycle.

\subsection{\texttt{forge} --- Scoped Construction}\label{subsec:forge}

\index{forge@\texttt{forge}|(}

A \lstinline|forge| block lets you construct a complex crystal value through a series of mutable operations, automatically freezing the result when the block completes.
Think of it like a foundry: the material is molten inside the forge, but what comes out is solid.

\begin{lstlisting}[language=Lattice, caption={Building a crystal value with \texttt{forge}}]
fix server_config = forge {
    flux temp = Map::new()
    temp.set("host", "localhost")
    temp.set("port", "8080")
    temp.set("max_connections", "1000")
    temp.set("timeout_ms", "30000")
    freeze(temp)
}

print(phase_of(server_config))  // "crystal"
print(server_config.get("host"))  // "localhost"
\end{lstlisting}

The \lstinline|forge| block is especially valuable when the construction process requires loops, conditionals, or function calls:

\begin{lstlisting}[language=Lattice, caption={Complex construction inside \texttt{forge}}]
fix lookup_table = forge {
    flux table = Map::new()
    for i in 0..100 {
        table.set(to_string(i), i * i)
    }
    freeze(table)
}

print(lookup_table.get("7"))   // 49
print(lookup_table.get("12"))  // 144
\end{lstlisting}

The phase checker knows that a \lstinline|forge| block always produces a crystal result.
This is reflected in \filename{src/phase\_check.c}, where \lstinline|EXPR_FORGE| returns \lstinline|PHASE_CRYSTAL|.

\index{forge@\texttt{forge}|)}

\subsection{\texttt{anneal} --- Transform a Crystal}\label{subsec:anneal}

\index{anneal@\texttt{anneal}|(}

In metallurgy, \emph{annealing} is the process of heating a metal, transforming it, and letting it cool back to a hardened state.
Lattice's \lstinline|anneal| does the same: it thaws a crystal value, applies a transformation closure, and re-freezes the result.

\begin{lstlisting}[language=Lattice, caption={Annealing a crystal value}]
fix prices = freeze([9.99, 14.99, 19.99, 24.99])

// Apply a 10% discount without manually thawing/re-freezing
fix discounted = anneal(prices, |items: any| {
    flux result = []
    for price in items {
        result.push(price * 0.9)
    }
    return result
})

print(discounted)             // [8.991, 13.491, 17.991, 22.491]
print(phase_of(discounted))   // "crystal"
\end{lstlisting}

\lstinline|anneal| requires its first argument to be a crystal value.
If you pass a fluid value, the runtime produces an error.
The transformation closure receives the thawed (fluid) copy, and whatever it returns is automatically frozen.

\index{anneal@\texttt{anneal}|)}

\subsection{\texttt{crystallize} and \texttt{borrow} --- Scoped Phase Changes}\label{subsec:crystallize-borrow}

\index{crystallize@\texttt{crystallize}|(}
\index{borrow@\texttt{borrow}|(}

Sometimes you want to temporarily change a value's phase within a specific scope.
\lstinline|crystallize| temporarily freezes a variable for the duration of a block, then restores its original phase:

\begin{lstlisting}[language=Lattice, caption={Temporary crystallization}]
flux data = [1, 2, 3]

crystallize(data) {
    // Inside this block, data is crystal
    print(phase_of(data))  // "crystal"
    // data.push(4)  // This would error!
}

// Back to fluid outside
print(phase_of(data))  // "fluid"
data.push(4)
print(data)  // [1, 2, 3, 4]
\end{lstlisting}

\lstinline|borrow| is the inverse: it temporarily thaws a crystal variable so you can modify it within a block, then re-freezes it:

\begin{lstlisting}[language=Lattice, caption={Temporarily borrowing a crystal value}]
fix config = freeze(Map::new())

borrow(config) {
    // Inside this block, config is fluid
    config.set("debug", "true")
    config.set("verbose", "true")
}

// config is crystal again outside the borrow block
print(phase_of(config))  // "crystal"
\end{lstlisting}

These scoped operations are compiled to a sequence of phase tag changes and are verified by the phase checker (see \filename{src/stackcompiler.c}).
The key property: both \lstinline|crystallize| and \lstinline|borrow| restore the original phase when the block ends, even if an error occurs during execution.

\index{crystallize@\texttt{crystallize}|)}
\index{borrow@\texttt{borrow}|)}

\subsection{\texttt{sublimate} --- The Point of No Return}\label{subsec:sublimate}

\index{sublimate@\texttt{sublimate}|(}
\index{sublimated phase}

In chemistry, sublimation is when a solid transforms directly into gas, skipping the liquid state entirely.
In Lattice, \lstinline|sublimate| converts a value into the \emph{sublimated} phase (\lstinline|VTAG_SUBLIMATED|)---a terminal, irreversible state of immutability.

\begin{lstlisting}[language=Lattice, caption={Sublimating a value}]
flux secret_key = "my-api-key-12345"
sublimate(secret_key)

print(phase_of(secret_key))  // "sublimated"
// secret_key cannot be thawed, reassigned, or modified
// It is permanently locked
\end{lstlisting}

Unlike crystal values, sublimated values \emph{cannot be thawed}.
Calling \lstinline|thaw()| on a sublimated value produces an error.
This makes sublimation ideal for security-sensitive data (API keys, passwords, certificates) or constants that must never change under any circumstances.

\begin{warnbox}{Sublimation Is Irreversible}
Once a value is sublimated, there is no way to make it mutable again within the runtime.
Use sublimation deliberately for values that truly must be permanent.
\end{warnbox}

\index{sublimate@\texttt{sublimate}|)}

\section{When and Why to Freeze Values}\label{sec:when-to-freeze}

\index{freeze@\texttt{freeze()}!when to use}

Knowing \emph{how} to freeze is only half the story.
Knowing \emph{when} to freeze is what separates a novice Lattice programmer from an expert.
Here are the primary scenarios where freezing pays dividends.

\subsection{Sharing Data Between Threads}

In Lattice's structured concurrency model (covered in \cref{ch:structured-concurrency}), spawned tasks should not share mutable state.
The strict mode phase checker enforces this: fluid bindings cannot cross a spawn boundary.
The solution is to freeze data before sharing:

\begin{lstlisting}[language=Lattice, caption={Freezing data for safe sharing across tasks}]
flux results = [10, 20, 30, 40, 50]
fix shared = freeze(results)

scope {
    spawn {
        // Safe: shared is crystal, no data races
        print("Task sees: " + to_string(shared))
    }
    spawn {
        print("Other task sees: " + to_string(shared))
    }
}
\end{lstlisting}

Crystal values are inherently safe to share because no one can modify them.
This eliminates an entire class of concurrency bugs---no locks needed, no synchronization overhead, no data races.

\subsection{API Boundaries}

When designing a function that returns data the caller should not modify, freeze it:

\begin{lstlisting}[language=Lattice, caption={Returning frozen data from a function}]
fn get_default_settings() {
    flux defaults = Map::new()
    defaults.set("theme", "dark")
    defaults.set("font_size", "14")
    defaults.set("language", "en")
    return freeze(defaults)
}

fix settings = get_default_settings()
// Callers get crystal data---they can read but not modify
\end{lstlisting}

\subsection{Snapshot Semantics}

When you need to capture the state of a value at a particular moment, freeze it:

\begin{lstlisting}[language=Lattice, caption={Creating snapshots with freeze}]
flux log_entries = []
log_entries.push("Server started")
log_entries.push("Listening on port 8080")

fix startup_log = freeze(log_entries)

log_entries.push("Connection from 192.168.1.1")
log_entries.push("Request: GET /api/status")

fix runtime_log = freeze(log_entries)

print(startup_log.len())  // 2
print(runtime_log.len())  // 4
\end{lstlisting}

Each \lstinline|freeze()| call captures the array at that point in time.
The snapshots are independent and immutable, while the original continues to grow.

\subsection{Performance Optimization}

Crystal values stored in arena-backed CrystalRegions enjoy better cache locality than fluid values scattered across the general heap.
For read-heavy workloads with large data structures, freezing the data can improve access performance.
We explore the memory implications in \cref{ch:memory-arenas}.

\subsection{Correctness Guarantees}

Sometimes you freeze a value simply to catch bugs.
If you know a value should not change after a certain point, freezing it converts silent corruption into a loud runtime error:

\begin{lstlisting}[language=Lattice, caption={Freezing for correctness}]
fn validate_and_seal(data: any) {
    // ... validation logic ...
    return freeze(data)
}

fix user_record = validate_and_seal(raw_input)
// Any code that accidentally tries to mutate user_record
// will get an immediate phase error instead of silent corruption
\end{lstlisting}

\section{Phase Errors and What They Tell You}\label{sec:phase-errors}

\index{phase errors|(}

When you violate the phase system's rules, Lattice produces clear error messages.
Understanding these errors is essential for working effectively with phases.

\subsection{Mutating a Crystal Value}

The most common phase error: trying to modify something that is frozen.

\begin{lstlisting}[language=Lattice, caption={Crystal mutation error}]
fix names = freeze(["Alice", "Bob"])
// names.push("Charlie")
// Error: cannot mutate crystal value
\end{lstlisting}

This error means you tried to call a mutating method (\lstinline|push|, \lstinline|set|, \lstinline|pop|, etc.) on a crystal value.
The fix is to either \lstinline|thaw()| the value first or restructure your code so the mutation happens before the freeze.

\subsection{Assigning to a Crystal Binding}

\begin{lstlisting}[language=Lattice, caption={Crystal assignment error}]
fix count = 42
// count = 43
// Error: cannot reassign fix binding 'count'
\end{lstlisting}

A \lstinline|fix| binding cannot be reassigned.
If you need a new value, create a new binding:

\begin{lstlisting}[language=Lattice, caption={Working with crystal bindings}]
fix count = 42
fix new_count = count + 1
print(new_count)  // 43
\end{lstlisting}

\subsection{Strict Mode: \texttt{let} Not Allowed}

\begin{lstlisting}[language=Lattice, caption={Strict mode rejects \texttt{let}}]
#mode strict

// let x = 10
// Error: strict mode: use 'flux' or 'fix' instead of 'let'
//        for binding 'x'

flux x = 10  // Use this instead
\end{lstlisting}

In strict mode, every binding must have an explicit phase.
This forces you to think about mutability at every declaration point.

\subsection{Strict Mode: Phase Mismatch}

\begin{lstlisting}[language=Lattice, caption={Strict mode phase mismatch}]
#mode strict

fix frozen_val = freeze(42)
// flux reassigned = frozen_val
// Error: cannot bind crystal value with flux for 'reassigned'
\end{lstlisting}

In strict mode, the phase checker tracks the phase of every expression and ensures that bindings are consistent.
Binding a crystal expression with \lstinline|flux| is an error because it would create a contradiction: the value is crystal but the binding says it should be fluid.

\subsection{Strict Mode: Freezing an Already Crystal Value}

\begin{lstlisting}[language=Lattice, caption={Double-freeze error in strict mode}]
#mode strict

fix data = freeze([1, 2, 3])
// fix again = freeze(data)
// Error: strict mode: cannot freeze an already crystal value
\end{lstlisting}

This error helps catch redundant operations.
If a value is already crystal, freezing it again is pointless---and in strict mode, Lattice tells you so.

\subsection{Strict Mode: Thawing an Already Fluid Value}

\begin{lstlisting}[language=Lattice, caption={Thawing a fluid value in strict mode}]
#mode strict

flux data = [1, 2, 3]
// flux copy = thaw(data)
// Error: strict mode: cannot thaw an already fluid value
\end{lstlisting}

Similarly, thawing a fluid value is redundant.
If you want an independent copy of a fluid value, use \lstinline|clone()| instead.

\subsection{Strict Mode: Fluid Values in \texttt{spawn}}

\begin{lstlisting}[language=Lattice, caption={Fluid value crossing spawn boundary}]
#mode strict

flux shared_counter = 0

scope {
    spawn {
        // print(shared_counter)
        // Error: strict mode: cannot use fluid binding
        //        'shared_counter' across thread boundary in spawn
    }
}
\end{lstlisting}

The phase checker in \filename{src/phase\_check.c} includes a special \lstinline|pc_check_spawn_expr()| function that walks expressions inside \lstinline|spawn| blocks and errors if any fluid binding is referenced.
This is one of the most valuable strict mode checks: it prevents data races at compile time.

\subsection{Annotation Violations}

\begin{lstlisting}[language=Lattice, caption={Phase annotation violation}]
#mode strict

@crystal
flux temperature = 72.5
// Error: @crystal annotation violated: initializer for
//        'temperature' is fluid
\end{lstlisting}

Phase annotations (\lstinline|@fluid|, \lstinline|@crystal|) are assertions about what phase a value should have.
If the initializer's phase contradicts the annotation, the phase checker produces an error.
We cover annotations in detail in \cref{ch:strict-mode}.

\begin{tipbox}{Reading Phase Errors}
Phase errors always tell you three things:
\begin{enumerate}
  \item \textbf{What went wrong}: the specific phase violation.
  \item \textbf{Where it happened}: the variable name or expression involved.
  \item \textbf{Why it matters}: strict mode errors explain the rule being enforced.
\end{enumerate}
When you see a phase error, ask yourself: ``Should this value be mutable here, or should I have frozen it earlier?''
The answer usually reveals the right fix.
\end{tipbox}

\index{phase errors|)}

\section{Querying Phase at Runtime}\label{sec:querying-phase}

\index{phase\_of@\texttt{phase\_of()}}

The built-in \lstinline|phase_of()| function returns the current phase of any value as a string:

\begin{lstlisting}[language=Lattice, caption={Using \texttt{phase\_of()} for runtime inspection}]
flux counter = 0
fix pi = 3.14159
let name = "Lattice"

print(phase_of(counter))  // "fluid"
print(phase_of(pi))       // "crystal"
print(phase_of(name))     // "unphased"
\end{lstlisting}

The four possible return values are:
\begin{itemize}
  \item \lstinline|"fluid"| --- the value is mutable (\lstinline|VTAG_FLUID|).
  \item \lstinline|"crystal"| --- the value is immutable (\lstinline|VTAG_CRYSTAL|).
  \item \lstinline|"unphased"| --- the value has no explicit phase (\lstinline|VTAG_UNPHASED|).
  \item \lstinline|"sublimated"| --- the value is permanently immutable (\lstinline|VTAG_SUBLIMATED|).
\end{itemize}

The \lstinline|phase_of()| function is implemented as a native built-in registered in \filename{src/runtime.c}.
It inspects the \lstinline|phase| field of the \lstinline|LatValue| struct directly, making it a zero-cost query.

\lstinline|phase_of()| is particularly useful in library code that needs to handle values of any phase:

\begin{lstlisting}[language=Lattice, caption={Phase-aware library function}]
fn safe_append(collection: any, item: any) {
    if phase_of(collection) == "crystal" {
        // Can't modify crystal, return a new thawed copy
        flux copy = thaw(collection)
        copy.push(item)
        return freeze(copy)
    }
    // Fluid: modify in place
    collection.push(item)
    return collection
}

flux items = [1, 2, 3]
safe_append(items, 4)
print(items)  // [1, 2, 3, 4]

fix frozen = freeze([10, 20])
fix updated = safe_append(frozen, 30)
print(updated)  // [10, 20, 30]
\end{lstlisting}

\section{Summary: The Phase Lifecycle}\label{sec:phase-lifecycle}

Let us put it all together.
A value's phase lifecycle in Lattice looks like this:

\begin{enumerate}
  \item A value is \textbf{created} via a constructor, literal, or expression. It begins as \lstinline|VTAG_UNPHASED|.
  \item When bound with \lstinline|flux|, it transitions to \lstinline|VTAG_FLUID| and lives on the FluidHeap under GC management.
  \item When bound with \lstinline|fix| or passed through \lstinline|freeze()|, it transitions to \lstinline|VTAG_CRYSTAL| and may be relocated to a CrystalRegion.
  \item \lstinline|thaw()| creates a \emph{new} fluid copy from a crystal value. \lstinline|clone()| creates a copy at the same phase.
  \item \lstinline|sublimate()| converts to the \lstinline|VTAG_SUBLIMATED| terminal state---no going back.
  \item \lstinline|forge|, \lstinline|anneal|, \lstinline|crystallize|, and \lstinline|borrow| provide scoped and structured phase transitions.
\end{enumerate}

\begin{figure}[h]
\centering
\begin{tikzpicture}[
    phase/.style={draw, rounded corners, minimum width=2.5cm, minimum height=1cm, font=\small\bfseries},
    fluid/.style={phase, fill=blue!15, draw=latticeblue},
    crystal/.style={phase, fill=purple!15, draw=latticepurple},
    unphased/.style={phase, fill=gray!15, draw=latticegray},
    sublimated/.style={phase, fill=orange!15, draw=latticeorange},
    arr/.style={-{Stealth[length=6pt]}, thick},
]
    \node[unphased] (unphased) at (0,0) {Unphased};
    \node[fluid] (fluid) at (-3,-2.5) {Fluid};
    \node[crystal] (crystal) at (3,-2.5) {Crystal};
    \node[sublimated] (sublimated) at (3,-5) {Sublimated};

    \draw[arr] (unphased) -- node[left, font=\small] {\texttt{flux}} (fluid);
    \draw[arr] (unphased) -- node[right, font=\small] {\texttt{fix}} (crystal);
    \draw[arr, bend left=15] (fluid) to node[above, font=\small] {\texttt{freeze()}} (crystal);
    \draw[arr, bend left=15] (crystal) to node[below, font=\small] {\texttt{thaw()}} (fluid);
    \draw[arr] (crystal) -- node[right, font=\small] {\texttt{sublimate()}} (sublimated);
    \draw[arr] (fluid) to[bend right=25] node[left, font=\small] {\texttt{sublimate()}} (sublimated);
\end{tikzpicture}
\caption{The phase transition diagram for Lattice values.}\label{fig:phase-transitions}
\end{figure}

\section{Exercises}\label{sec:ch11-exercises}

\begin{enumerate}
  \item \textbf{Phase Explorer.}
    Write a function \lstinline|describe_phase(val: any)| that takes any value and prints a message describing its phase.
    Include different messages for fluid, crystal, unphased, and sublimated values.
    Test it with values of each phase.

  \item \textbf{Immutable Stack.}
    Implement a stack data structure using frozen arrays.
    Your \lstinline|push| operation should return a \emph{new} frozen array with the element added, and your \lstinline|pop| operation should return a tuple of the top element and the remaining frozen array.
    The original stack should never change.

  \item \textbf{Configuration Builder.}
    Write a \lstinline|build_config()| function that uses a \lstinline|forge| block to construct a configuration map.
    The function should accept an array of key-value tuples and build the frozen map from them.

  \item \textbf{Phase-Safe Merge.}
    Write a function \lstinline|merge_arrays(a: any, b: any)| that concatenates two arrays regardless of their phases.
    If either input is crystal, the function should thaw it before merging.
    The result should always be frozen.

  \item \textbf{Error Collector.}
    In strict mode, write a program that deliberately triggers each of the phase errors described in \cref{sec:phase-errors} (one at a time).
    Comment out each error after confirming the message, and write a brief comment explaining what the error means.
\end{enumerate}

\bigskip

\noindent\textbf{What's Next.}\quad
Now that we understand the phase system's semantics, we are ready to peer beneath the surface and ask: \emph{where do these values actually live?}
In \cref{ch:memory-arenas}, we explore the dual-heap memory architecture that makes the phase system possible---the FluidHeap for mutable values, the CrystalRegion for frozen data, and the ephemeral BumpArena for short-lived temporaries.
Understanding this architecture will give you the intuition to write code that is not only correct but efficient.

\index{phase system|)}
